\begin{pastexam}[label={px:nccu-109-dsa-01}]{政大資科所 109 學年度 第 1 題}
（問答）

(20 points) What are the worst case time complexities of the following algorithms?
\begin{enumerate}[label=1.\arabic*]
\item Binary search in a sorted array consisting of $n$ integers
\item Building a binary heap from an unsorted array consisting of $n$ integers
\item Randomized quicksort for an array consisting of $n$ integers
\item Kruskal algorithm with disjoint-set-forest for finding the minimum spanning tree given a
connected, undirected graph with $|V|$ vertices and $|E|$ edges
\end{enumerate}
\end{pastexam}

\begin{pastsoln}
答案：
\begin{center}
\begin{tabular}{cccc}
\toprule
1.1 & 1.2 & 1.3 & 1.4 \\
\midrule
$\Theta(\log n)$ & $\Theta(n)$ & $\Theta(n^2)$ & $O(|E| \log |V|) = O(|E| \log |E|)$ \\
\bottomrule
\end{tabular}
\end{center}

\textbf{1.1 二分搜尋}（資結書第~17.1~節）：每比一次，剩下的範圍至少減半，長度 $n$ 的範圍最多比
$\lfloor \lg n \rfloor + 1$ 次就變空，最壞（要找的數不在陣列裡）比這麼多次，$\Theta(\log n)$。

\textbf{1.2 建二元堆積}（資結書第~8.2~節）：由下而上，對 $i = \lfloor n/2 \rfloor, \ldots, 1$ 做下沉。
節點 $i$ 下沉的成本和它的高度成正比，所有節點的高度加總小於 $2n$（資結書定理~8.1），所以最壞是
$\Theta(n)$，不是直覺的 $O(n \log n)$。 如果改用「一個一個插入」建堆積，最壞是 $\Theta(n \log n)$
（例如插入遞增的數到最大堆積，每個都要上浮到根），題目問的是從陣列建堆積，一般指前者。

\textbf{1.3 隨機化快速排序}（資結書第~14.5~節）：\textbf{期望}時間是 $\Theta(n \log n)$，但題目問最壞情況。
隨機選基準仍然可能（機率極小）每次都選到範圍裡的最大或最小值，每次只切掉一個數，
$T(n) = T(n-1) + \Theta(n) = \Theta(n^2)$。 任兩個數最多比較一次（資結書引理~14.1），所以也不會超過
$O(n^2)$。 最壞是 $\Theta(n^2)$。

\textbf{1.4 Kruskal}（資結書第~24.2~節）：
\begin{itemize}
\item 邊依權重排序：$O(|E| \log |E|)$。
\item $|V|$ 次 \proc{Make-Set}、$2|E|$ 次 \proc{Find-Set}、$|V| - 1$ 次 \proc{Union}：用按秩合併加路徑壓縮
的互斥集合森林，總共 $O((|V| + |E|)\,\alpha(|V|))$（資結書定理~13.4），$\alpha$ 增長極慢，比排序小。
\end{itemize}
總共 $O(|E| \log |E|)$。 圖是連通的，$|V| - 1 \le |E| \le |V|^2$，所以 $\log |E| \le 2 \log |V|$，
可以寫成 $O(|E| \log |V|)$。 排序的部分是比較排序的最壞下界，所以這也是緊的界（權重已經排好或是小整數
時才能更快）。 只用按秩合併、不做路徑壓縮時每次 \proc{Find-Set} 是 $O(\log |V|)$，總時間一樣是
$O(|E| \log |V|)$。
\end{pastsoln}
